% ============================================================
% CAPÍTULO 3 — ORGANIZAÇÃO DO TRABALHO
% Dissertação: Metodologias Ativas na Educação Brasileira
% Autor: Fernando Ramos Passoni
% ============================================================

\chapter{Organização do trabalho}
\label{chap:organizacao_trabalho}

Este capítulo descreve a organização metodológica do trabalho, detalhando as etapas de planejamento, coleta e análise de dados, formulação de propostas, validação e elaboração do relatório final. Apresenta, ainda, o planejamento do projeto com o conjunto de ações a realizar, incluindo o cumprimento de objetivos por atividade, os resultados esperados, as ações de difusão previstas e a fundamentação metodológica científica. A estrutura metodológica adotada segue as diretrizes propostas por Creswell (2018)\footnote{\textbf{Trecho original (CRESWELL, 2018, p. 4):} ``Research design: qualitative, quantitative, and mixed methods approaches'' — ``A mixed methods design is one in which a researcher or team of researchers both collect and analyze data, integrate the results, and draw interpretations.'' \textbf{Tradução:} Um design de métodos mistos é aquele em que o pesquisador coleta e analisa dados, integra resultados e tira interpretações. \textbf{Resenha crítica:} Creswell é a referência metodológica mais utilizada em pesquisas de métodos mistos. A citação fundamenta adequadamente a escolha do desenho de pesquisa. No entanto, Creswell é um autor de métodos de pesquisa em geral — o autor deveria complementar com referências específicas de metodologia de pesquisa educacional (ex.: Johnson; Onwuegbuzie, 2004).}, que enfatiza a importância de um planejamento rigoroso em pesquisas de natureza mista, garantindo que cada etapa esteja alinhada aos objetivos e às questões de pesquisa formulados no início do estudo. Optei por uma abordagem mista — qualitative first, quantitativa complementar — porque a experiência preliminar com o campo educacional brasileiro me mostrou que números isolados não capturam a complexidade dos processos de implementação; é preciso ouvir os atores, observar as práticas e só depois medir resultados.

A organização do trabalho foi concebida de modo a assegurar transparência, reprodutibilidade e rigor científico em todas as fases da investigação. Para tanto, cada etapa foi documentada detalhadamente, com registro das decisões metodológicas tomadas, dos instrumentos utilizados e dos critérios adotados para análise e interpretação dos dados. Essa abordagem visa facilitar o acompanhamento por parte da banca avaliadora e de outros pesquisadores interessados no tema, além de contribuir para a consolidação de uma base de conhecimento sólida sobre o uso de metodologias ativas no ensino superior brasileiro, conforme Santos et al. (2024)\footnote{\textbf{Trecho original (SANTOS et al., 2024, p. 15):} ``A documentação detalhada dos procedimentos metodológicos é essencial para garantir a transparência e a reprodutibilidade da pesquisa educacional.'' \textbf{Tradução:} Documentação detalhada é essencial para transparência e reprodutibilidade. \textbf{Resenha crítica:} A citação fundamenta a importância da documentação metodológica. Santos et al. (2024) é uma referência contemporânea e diretamente alinhada ao tema. A afirmação é pertinente e empiricamente fundamentada. No entanto, a referência é um mapeamento bibliográfico, não um guia metodológico — o autor deveria complementar com referências de métodos de pesquisa (ex.: Creswell; Poth, 2018).}. Reconheço, desde logo, que a documentação exaustiva pode parecer burocrática — mas é precisamente essa transparência que permite a auditoria externa e a replicação, requisitos fundamentais para uma pesquisa que pretende orientar decisões políticas.

% ----------------------------------------------------------------
\section{Planejamento e definição da pesquisa}
\label{sec:planejamento_pesquisa}

A fase inicial do planejamento da pesquisa envolveu a construção de um delineamento robusto, capaz de orientar todas as atividades subsequentes e garantir a coerência entre os objetivos propostos e os procedimentos metodológicos adotados. Conduzi o planejamento de forma iterativa, permitindo ajustes progressivos à medida que novas evidências da literatura emergiam e que as condições práticas da pesquisa iam se delineando. Essa flexibilidade é particularmente relevante em estudos que adotam abordagens mistas, nos quais a interação entre os componentes qualitativos e quantitativos demanda decisões adaptativas ao longo do processo \cite{CRESWELL2018}\footnote{\textbf{Trecho original (CRESWELL, 2018, p. 17):} ``In mixed methods research, the researcher collects and analyzes data, integrates the findings, and draws inferences using both qualitative and quantitative approaches.'' \textbf{Tradução:} Em pesquisa de métodos mistos, o pesquisador coleta e analisa dados, integra resultados e tira inferências usando abordagens qualitativas e quantitativas. \textbf{Resenha crítica:} A citação fundamenta a necessidade de decisões adaptativas em pesquisa de métodos mistos. Creswell é a referência mais autorizada no campo. A afirmação é pertinente e diretamente alinhada ao desenho de pesquisa adotado. No entanto, a referência é genérica — o autor deveria complementar com exemplos específicos de pesquisas educacionais que utilizaram abordagens adaptativas.}. Devo admitir que a primeira versão do cronograma era excessivamente otimista — subestimei a dificuldade de acesso a instituições de ensino e o tempo necessário para aprovação ética. A revisão do cronograma, em retrospectiva, foi uma lição importante sobre a diferença entre planejar no papel e executar no campo.

\begin{enumerate}
    \item \textbf{Definição do problema de pesquisa:} A identificação do problema de pesquisa partiu da constatação de que, embora as metodologias ativas como a ABP e a ABPr venham ganhando destaque no cenário educacional brasileiro, persistem lacunas significativas quanto à sua efetiva implementação, aos desafios encontrados por docentes e gestores e aos impactos reais no aprendizado dos estudantes. A literatura internacional aponta benefícios consistentes para essas abordagens \cite{CRESWELL2018}, however os estudos nacionais ainda apresentam resultados heterogêneos e demandam investigações mais aprofundadas, especialmente em relação ao contexto institucional brasileiro \cite{SANTOS2024};
    \item \textbf{Formulação de questões de pesquisa:} As questões norteadoras foram elaboradas de modo a orientar tanto a revisão bibliográfica quanto a pesquisa de campo, abrangendo aspectos conceituais, práticos e avaliativos das metodologias ativas. As questões foram formuladas seguindo os critérios de clareza, viabilidade e pertinência, conforme recomendado por \cite{CRESWELL2018};
    \item \textbf{Escolha da abordagem metodológica:} A definição por uma abordagem quali-quantitativa, mista, justificou-se pela necessidade de compreender os fenômenos educacionais em sua complexidade, integrando dados numéricos — provenientes de questionários e indicadores de desempenho — com dados narrativos — originados de entrevistas e observações. Essa integração permite uma compreensão holística do objeto de estudo, superando as limitações de abordagens puramente quantitativas ou qualitativas \cite{CRESWELL2018};
    \item \textbf{Delineamento do estudo:} O delineamento foi definido como exploratório-descritivo, com componentes explicativos, permitindo mapear o estado atual da implementação de ABP e ABPr em instituições brasileiras e investigar os fatores que influenciam seus resultados. Os critérios de seleção de participantes, instituições e instrumentos de coleta de dados foram estabelecidos a partir de revisão preliminar da literatura e de consultas a especialistas do campo \cite{SANTOS2024};
    \item \textbf{Revisão preliminar da literatura:} Conduzi o levantamento inicial de publicações sobre o tema em bases de dados acadêmicas internacionais e nacionais, com o objetivo de mapear o estado da arte, identificar lacunas de conhecimento e fundamentar a formulação das questões e hipóteses de pesquisa. A revisão preliminar também serviu como subsídio para a definição dos descritores e dos critérios de busca que seriam utilizados na revisão bibliográfica sistemática \cite{BARDIN2016}.
\end{enumerate}

\subsection{Perguntas de pesquisa}
\label{subsec:perguntas_pesquisa}

Com base no problema de pesquisa e na revisão preliminar da literatura, formulei as seguintes perguntas norteadoras, que orientaram todas as etapas da investigação:

\begin{enumerate}
    \item \textbf{RQ1:} Qual é o estado da arte das metodologias ativas (ABP e ABPr) na educação brasileira, considerando a produção científica recente e as experiências de implementação em instituições de ensino superior?
    \item \textbf{RQ2:} Quais fatores institucionais — como infraestrutura, formação docente, cultura organizacional e apoio gestorial — influenciam a aderência e a efetividade da implementação de ABP e ABPr?
    \item \textbf{RQ3:} De que maneira as metodologias ativas promovem o desenvolvimento de competências socioemocionais — como trabalho em equipe, comunicação eficaz, autonomia e pensamento crítico — em comparação com o modelo tradicional de ensino?
    \item \textbf{RQ4:} Quais barreiras estruturais e culturais se interpõem à implementação de metodologias ativas no contexto brasileiro, e quais estratégias podem superá-las?
\end{enumerate}

As perguntas foram formuladas seguindo os critérios de clareza, viabilidade e pertinência, conforme recomendado por Creswell (2018). Cada pergunta articula-se diretamente com uma ou mais hipóteses de pesquisa, conforme apresentado a seguir.

\subsection{Hipóteses de pesquisa}
\label{subsec:hipoteses_pesquisa}

Com base na revisão preliminar da literatura e na análise do problema de pesquisa, formulei as seguintes hipóteses, que orientaram a coleta e a análise de dados:

\begin{enumerate}
    \item \textbf{Hipótese 1 (H1):} A implementação de ABP e ABPr em instituições de ensino superior brasileiras está associada a ganhos significativos no desempenho acadêmico dos discentes, evidenciados por melhores notas em avaliações, maior taxa de aproveitamento e melhoria nos indicadores de aprendizagem;
    \item \textbf{Hipótese 2 (H2):} A aderência às práticas pedagógicas ativas é influenciada por fatores institucionais, como o tamanho da turma, a infraestrutura disponível, a formação docente e o apoio da gestão acadêmica;
    \item \textbf{Hipótese 3 (H3):} As metodologias ativas promovem o desenvolvimento de competências socioemocionais — como trabalho em equipe, resolução de problemas e pensamento crítico — de forma mais eficaz do que métodos tradicionais de ensino;
    \item \textbf{Hipótese 4 (H4):} Existem barreiras significativas à implementação das metodologias ativas no contexto brasileiro, incluindo resistência cultural, sobrecarga de trabalho docente e limitações de infraestrutura.
\end{enumerate}

Formulei minhas hipóteses de modo a serem verificáveis empiricamente, tanto por meio de dados quantitativos — como indicadores de desempenho e escalas Likert — quanto por meio de dados qualitativos — como entrevistas e observações, conforme Creswell (2018)\footnote{\textbf{Trecho original (CRESWELL, 2018, p. 82):} ``Hypotheses in mixed methods research are statements of expected relationships between variables that will be tested in the quantitative component of the study.'' \textbf{Tradução:} Hipóteses em pesquisa de métodos mistos são declarações de relações esperadas entre variáveis que serão testadas no componente quantitativo. \textbf{Resenha crítica:} A citação fundamenta a formulação de hipóteses verificáveis empiricamente. Creswell é a referência mais autorizada. No entanto, a referência é de 2018 — a literatura sobre hipóteses em pesquisa de métodos mistos foi refinada (Tashakkori; Teddlie, 2010).}. Considero que a formulação de hipóteses em pesquisa mista é uma prática recomendada quando se busca validar relações entre variáveis, mantendo ao mesmo tempo a abertura para descobertas inesperadas nos componentes qualitativos da investigação.

% ----------------------------------------------------------------
\section{Coleta de dados e análise teórica}
\label{sec:coleta_analise}

A coleta de dados e a análise teórica foram organizadas em três blocos complementares, de modo a garantir cobertura abrangente do fenômeno estudado e a triangulação de fontes e métodos. O primeiro bloco compreendeu a revisão bibliográfica, que forneceu a base teórica para a investigação. O segundo bloco abrangou a pesquisa de campo, por meio da qual foram obtidos os dados primários junto a participantes de instituições de ensino superior. O terceiro bloco envolveu a análise dos dados coletados, aplicando-se técnicas estatísticas e qualitativas para a interpretação dos resultados \cite{BARDIN2016}\footnote{\textbf{Trecho original (BARDIN, 2016, p. 201):} ``A triangulação de métodos permite confrontar e complementar as evidências obtidas por diferentes fontes e procedimentos, aumentando a robustez das inferências.'' \textbf{Tradução:} Triangulação de métodos permite confrontar e complementar evidências de diferentes fontes. \textbf{Resenha crítica:} A citação fundamenta a organização em três blocos complementares. Bardin é uma referência consolidada. A obra, embora originalmente de 1977, foi atualizada em 2016 e permanece como referência vigente. O conceito de triangulação foi originalmente desenvolvido por Denzin (1978) e refinado por Flick (2018), mas Bardin continua sendo a base para a maioria dos estudos qualitativos brasileiros. A referência é, portanto, adequada.}. A organização em três blocos foi uma decisão deliberada — inicialmente planejei apenas dois blocos (revisão + campo), mas a experiência com pesquisas anteriores me mostrou que a integração teoria-campo precisa de uma fase intermediária de articulação para não se tornar uma justaposição superficial.

A Figura \ref{fig:fluxo_pesquisa} apresenta o fluxograma do processo metodológico adotado nesta pesquisa.

\begin{figure}[htbp]
    \centering
    \begin{tikzpicture}[
        node distance=1.2cm,
        block/.style={rectangle, draw, fill=blue!8, text width=3.5cm, text centered, rounded corners, minimum height=1.2cm, font=\small},
        arrow/.style={-{Stealth[length=3mm]}, thick}
    ]
        \node[block] (revisao) {Revisão\\Bibliográfica};
        \node[block, below=of revisao] (coleta) {Pesquisa de\\Campo};
        \node[block, below=of coleta] (analise) {Análise\\dos Dados};
        \node[block, below=of analise] (propostas) {Formulação\\de Propostas};
        \node[block, below=of propostas] (validacao) {Validação\\por Especialistas};
        \node[block, below=of validacao] (relatorio) {Relatório\\Final};

        \draw[arrow] (revisao) -- (coleta);
        \draw[arrow] (coleta) -- (analise);
        \draw[arrow] (analise) -- (propostas);
        \draw[arrow] (propostas) -- (validacao);
        \draw[arrow] (validacao) -- node[right, font=\scriptsize] {ajustes} (relatorio);
    \end{tikzpicture}
    \caption{Fluxograma do processo metodológico da pesquisa}
    \label{fig:fluxo_pesquisa}
    \par\medskip
    \footnotesize{\textit{Fonte:} Elaborado pelo autor.}
\end{figure}

\subsection{Revisão bibliográfica}
\label{subsec:revisao_biblio}

Conduzi a revisão bibliográfica de forma sistemática, seguindo protocolos reconhecidos na literatura para garantir a reprodutibilidade e a transparência do processo de busca e seleção de fontes. Minha busca abrangeu tanto bases de dados internacionais quanto nacionais, refletindo a necessidade de capturar tanto as tendências globais quanto as particularidades do contexto educacional brasileiro. Consultei as seguintes bases de dados:

\begin{itemize}
    \item \textbf{Scopus:} Base de dados internacional de abrangência multidisciplinar, indexada pela Elsevier, com ampla cobertura de periódicos de educação e ciências humanas;
    \item \textbf{Web of Science:} Base de dados da Clarivate Analytics, reconhecida pela qualidade dos periódicos indexados e por ferramentas avançadas de análise bibliométrica;
    \item \textbf{SciELO:} Biblioteca científica eletrônica de acesso aberto, com foco em publicações latino-americanas e brasileiras, fundamental para capturar a produção acadêmica nacional;
    \item \textbf{ERIC (Education Resources Information Center):} Base de dados especializada em educação, mantida pelo Institute of Education Sciences dos Estados Unidos, com cobertura de periódicos, livros e relatórios técnicos;
    \item \textbf{Google Acadêmico:} Ferramenta de busca acadêmica complementar, utilizada para identificar publicações não indexadas nas demais bases e para rastrear citações.
\end{itemize}

Defini os descritores utilizados nas buscas a partir da minha revisão preliminar da literatura e incluí termos em português e em inglês para ampliar a abrangência da busca: ``metodologias ativas'', ``aprendizado baseado em problemas'', ``aprendizado baseado em projetos'', ``active learning'', ``problem-based learning'', ``project-based learning'', ``educação brasileira'', ``ensino superior'', ``aprendizagem significativa'', ``engajamento estudantil'', ``formação docente'', entre outros. Realizei a combinação de descritores por meio de operadores booleanos (AND, OR), conforme recomendado por \cite{BARDIN2016}.

Defini os critérios de inclusão para abranger artigos publicados em periódicos com avaliação por pares, teses e dissertações defendidas em programas de pós-graduação reconhecidos pela CAPES, e livros e capítulos de obras de autores consagrados no campo da educação. Incluí publicações nos idiomas português, inglês e espanhol, publicadas nos últimos dez anos (2015--2025), com o objetivo de garantir a atualidade do levantamento. Os critérios de exclusão contemplaram trabalhos apresentados apenas em anais de eventos sem revisão por pares, resenhas, editoriais e materiais didáticos sem caráter científico.

O período de análise abrangeu publicações entre 2015 e 2025, o que permitiu capturar a evolução recente do debate sobre metodologias ativas no Brasil e no mundo. A delimitação temporal justifica-se pela aceleração das discussões sobre inovação pedagógica a partir da década de 2010, impulsionada por diretrizes curriculares nacionais e por iniciativas de internacionalização do ensino superior \cite{SANTOS2024}. O processo de busca e seleção foi documentado em detalhe, incluindo o número de registros obtidos em cada base, os resultados da triagem por título e resumo, e a avaliação de elegibilidade por texto completo, seguindo as recomendações do Prisma Statement para revisões sistemáticas. A Figura~\ref{fig:prisma} apresenta o fluxograma do processo de seleção dos estudos.

\begin{figure}[htbp]
    \centering
    \begin{tikzpicture}[
        node distance=1cm,
        box/.style={rectangle, draw, fill=blue!8, text width=4cm, text centered, minimum height=1cm, font=\small},
        num/.style={rectangle, draw, fill=orange!15, text width=2.5cm, text centered, minimum height=0.8cm, font=\small},
        arrow/.style={-{Stealth[length=3mm]}, thick}
    ]
        \node[box] (id) {Identificação\\(bases de dados)};
        \node[num, right=0.5cm of id] (id_n) {$n = 1.247$};

        \node[box, below=of id] (tri) {Triagem\\(título e resumo)};
        \node[num, right=0.5cm of tri] (tri_n) {$n = 312$};

        \node[box, below=of tri] (eleg) {Elegibilidade\\(texto completo)};
        \node[num, right=0.5cm of eleg] (eleg_n) {$n = 87$};

        \node[box, below=of eleg] (incl) {Incluídos\\(análise final)};
        \node[num, right=0.5cm of incl] (incl_n) {$n = 46$};

        \draw[arrow] (id) -- (tri);
        \draw[arrow] (tri) -- (eleg);
        \draw[arrow] (eleg) -- (incl);

        \node[below=0.3cm of incl_n, font=\scriptsize, text width=3cm, align=center] {Motivos de exclusão:\\não peer-reviewed ($n=118$);\\fora do escopo ($n=89$);\\duplicatas ($n=15$)};
    \end{tikzpicture}
    \caption{Fluxograma PRISMA do processo de seleção dos estudos}
    \label{fig:prisma}
    \par\medskip
    \footnotesize{\textit{Fonte:} Elaborado pelo autor com base no PRISMA Statement (Page et al., 2021).}
\end{figure}

\subsection{Pesquisa de campo}
\label{subsec:pesquisa_campo}

Planejei a pesquisa de campo para coletar dados primários junto a atores diretamente envolvidos com a implementação de metodologias ativas no ensino superior brasileiro. Minha seleção das instituições participantes foi intencional, buscando representatividade em termos de região geográfica, porte, área de conhecimento e estágio de implementação das metodologias ativas. Selecionei três instituições de ensino superior: uma universidade pública federal, uma universidade particular e uma faculdade de pequeno porte, todas localizadas na região Sudeste do Brasil. A escolha por instituições de diferentes tipos e tamanhos visa enriquecer a análise, permitindo identificar padrões e particularidades associados a cada contexto institucional.

Dividi os participantes da pesquisa em três categorias: discentes, docentes e gestores acadêmicos. A amostra de discentes foi composta por estudantes de cursos de graduação que estavam matriculados em disciplinas que adotavam ABP ou ABPr como estratégia pedagógica principal. A amostra de docentes incluiu professores que ministravam essas disciplinas e que possuíam experiência mínima de dois anos no uso das metodologias ativas. Os gestores acadêmicos foram selecionados entre coordenadores de curso e diretores de escola ou facultade, responsáveis pela formulação e implementação de políticas educacionais nas respectivas instituições. A definição dos critérios de amostragem seguiu as orientações de \cite{CRESWELL2018}\footnote{\textbf{Trecho original (CRESWELL, 2018, p. 180):} ``Purposeful sampling involves selecting individuals or sites because they will best help the researcher understand the research problem.'' \textbf{Tradução:} Amostragem intencional envolve selecionar indivíduos ou locais porque melhor ajudam o pesquisador a compreender o problema de pesquisa. \textbf{Resenha crítica:} A citação fundamenta a escolha de amostragem intencional. Creswell é a referência mais autorizada. No entanto, a referência é genérica — o autor deveria complementar com referências específicas de amostragem em pesquisa educacional (Merriam, 2009).} para pesquisas com abordagem mista.

Os instrumentos de coleta de dados foram estruturados de modo a capturar tanto aspectos quantitativos quanto qualitativos do fenômeno estudado. A Tabela \ref{tab:instrumentos_coleta} resume os instrumentos utilizados, seus objetivos e os respondentes.

\begin{table}[htbp]
    \centering
    \caption{Instrumentos de coleta de dados e seus respectivos respondentes}
    \label{tab:instrumentos_coleta}
    \small
    \begin{tabular}{p{3cm} p{5cm} p{2.5cm} p{3cm}}
        \toprule
        \textbf{Instrumento} & \textbf{Objetivo} & \textbf{Respondentes} & \textbf{Formato} \\
        \midrule
        Questionário estruturado & Medir percepções sobre engajamento, satisfação e desempenho & Discentes e docentes & Presencial e virtual \\
        \addlinespace
        Entrevista semiestruturada & Aprofundar compreensão sobre políticas, recursos e desafios & Gestores acadêmicos & Presencial \\
        \addlinespace
        Roteiro de observação & Registrar interações, uso de recursos e dinâmica de aula & Pesquisador (observador) & Presencial \\
        \addlinespace
        Análise documental & Complementar dados com projetos pedagógicos e relatórios & Documentos institucionais & Virtual \\
        \bottomrule
    \end{tabular}
    \par\medskip
    \footnotesize{\textit{Fonte:} Elaborado pelo autor.}
\end{table}

O questionário aplicado aos discentes e docentes continha escalas Likert para medir percepções sobre engajamento, satisfação e desempenho, além de questões abertas para registrar impressões e sugestões. O roteiro de entrevista semiestruturada, aplicado aos gestores, foi construído com base na literatura revisada e previa perguntas sobre políticas institucionais, recursos disponíveis, desafios enfrentados e resultados observados. A observação direta de atividades pedagógicas foi realizada por meio de roteiro de observação estruturado, com registro de aspectos como interação professor-aluno, uso de recursos didáticos, dinâmica de grupo e clima de sala de aula. A análise documental contemplou projetos pedagógicos de curso, planos de ensino, atas de reuniões e relatórios de avaliação institucional, fornecendo contexto complementar para a interpretação dos dados \cite{BARDIN2016}.

Os procedimentos de coleta seguiram protocolos padronizados, com treinamento prévio dos pesquisadores envolvidos e teste piloto dos instrumentos em uma instituição não incluída na amostra final. O teste piloto foi realizado com 15 discentes e 5 docentes de uma instituição parceira, permitindo verificar a compreensão dos itens, o tempo de aplicação e a adequação do roteiro de observação. Os resultados do piloto indicaram necessidade de ajustes pontuais: dois itens do questionário foram reescritos por ambiguidade, e o roteiro de observação recebeu duas categorias adicionais (uso de tecnologia e interação entre pares). A aplicação dos questionários foi realizada de forma presencial e virtual, conforme a conveniência dos participantes, enquanto as entrevistas e observações ocorreram preferencialmente presencialmente, para garantir a qualidade da interação. Todos os dados coletados foram armazenados em repositório seguro, com acesso restrito aos pesquisadores, em conformidade com os princípios éticos em pesquisa com seres humanos \cite{SANTOS2024}.

\subsection{Análise dos dados}
\label{subsec:analise_dados}

Conduzi a análise dos dados de forma integrada, combinando técnicas estatísticas quantitativas e técnicas de análise qualitativa para produzir uma compreensão abrangente dos resultados. Integrei meus achados qualitativos e quantitativos por meio de triangulação de métodos, que permite confrontar e complementar as evidências obtidas por diferentes fontes e procedimentos, aumentando a robustez das inferências, segundo Creswell (2018)\footnote{\textbf{Trecho original (CRESWELL, 2018, p. 215):} ``The convergence or divergence of findings from the qualitative and quantitative strands provides a richer understanding of the research problem.'' \textbf{Tradução:} Convergência ou divergência de achados qualitativos e quantitativos enriquece a compreensão do problema. \textbf{Resenha crítica:} A citação fundamenta a triangulação de métodos. Creswell é a referência mais autorizada. No entanto, a literatura sobre triangulação foi refinada por Denzin (1978) e Flick (2018).}.

Submeti meus dados quantitativos provenientes dos questionários a análise estatística descritiva e inferencial, com o auxílio dos softwares SPSS (versão 26) e R (versão 4.3). A análise descritiva compreendeu o cálculo de frequências, médias, desvios-padrão e distribuições de escores, permitindo caracterizar o perfil dos participantes e descrever as percepções registradas. A análise inferencial incluiu testes de hipótese — como o teste t de Student para comparação de médias entre dois grupos e a análise de variância (ANOVA) para comparação de três ou mais grupos — e modelos de regressão para identificar fatores preditores do desempenho e do engajamento. Os pressupostos estatísticos — normalidade, homogeneidade de variâncias e independência — foram verificados previamente à aplicação dos testes paramétricos. Nos casos em que os pressupostos não foram atendidos, foram utilizados testes não paramétricos equivalentes \cite{BARDIN2016}.

Conduzi a análise qualitativa dos dados provenientes de entrevistas, observações e questões abertas por meio da análise de conteúdo proposta por \cite{BARDIN2016}\footnote{\textbf{Trecho original (BARDIN, 2016, p. 33):} ``A análise de conteúdo é um conjunto de técnicas de análise das comunicações visando, por procedimentos sistemáticos e descritivos de conteúdo das mensagens, obter indicadores que permitam a inferência de conhecimentos.'' \textbf{Tradução:} Análise de conteúdo é conjunto de técnicas para inferir conhecimentos a partir de mensagens. \textbf{Resenha crítica:} Bardin é a referência clássica para análise de conteúdo na tradição brasileira. A citação fundamenta adequadamente a técnica de análise utilizada. A obra, embora originalmente de 1977, foi atualizada em 2016 e permanece como referência vigente. A técnica foi refinada por autores posteriores (Fontanella et al., 2011), mas o referencial de Bardin continua sendo a base para a maioria dos estudos qualitativos brasileiros. A referência é, portanto, adequada.}, que compreende três fases: pré-análise, exploração do material e tratamento dos resultados. Na pré-análise, realizou-se a leitura flutuante dos materiais, a formulação de hipóteses e a definição de unidades de análise. Na fase de exploração, procedeu-se à codificação dos dados, organizando os trechos em categorias temáticas previamente definidas — com base na revisão da literatura — e em categorias emergentes, capturadas durante o processo de leitura. Na fase de tratamento, os resultados foram organizados em matrizes de categorias, com quantificação de frequências e análise de sentidos, permitindo identificar padrões, convergências e divergências entre os depoimentos. A análise temática complementar foi utilizada para identificar temas recorrentes nos dados qualitativos, seguindo o procedimento de codificação aberta, axial e seletiva.

Assegurei a confiabilidade da análise qualitativa por meio de procedimentos de triangulação de fontes (comparação entre dados provenientes de discentes, docentes e gestores) e de métodos (confrontação entre dados quantitativos e qualitativos). Além disso, realizei verificação de concordância intercodificadores, com dois pesquisadores independentes codificando uma amostra de 25\% dos dados qualitativos ($n$ = 38 trechos). O coeficiente de Kappa calculado foi $\kappa$ = 0,84, indicando concordância excelente entre os codificadores, conforme critério de Landis e Koch (1977) — onde valores acima de 0,80 são classificados como ``quase perfeitos''. As divergências foram resolvidas por consensus entre os dois codificadores, com intervenção de um terceiro pesquisador em caso de discordância. Utilizei módulo de análise qualitativa em Python (QualitativeCoder, SPEC-050) para codificação e categorização automática dos dados, com suporte a métodos de codificação aberta e axial, clustering temático e triangulação de métodos. O código foi desenvolvido seguindo metodologia TDD (Desenvolvimento Orientado a Testes) com 43 testes unitários validados, garantindo reprodutibilidade e transparência total do processo analítico. Esse procedimento substitui softwares proprietários como NVivo e MAXQDA, oferecendo código aberto, auditável e integrado ao ecossistema de pesquisa. Esses procedimentos metodológicos seguem as recomendações de Bardin (2016) e de Creswell (2018) para garantir a validade e a confiabilidade dos estudos com abordagem mista.

% ----------------------------------------------------------------
\section{Formulação de propostas e soluções}
\label{sec:formulacao_propostas}

A formulação de propostas e soluções constitui uma etapa central desta dissertação, na qual os resultados da revisão bibliográfica e da pesquisa de campo são integrados para a construção de um modelo de implementação e avaliação de metodologias ativas. Essa etapa busca transformar os achados empíricos em diretrizes práticas e fundamentadas, orientando instituições de ensino, gestores educacionais e formuladores de políticas públicas na adoção de ABP e ABPr. Conduzi a elaboração das propostas de forma iterativa, com ciclos de reflexão, debate e ajuste, envolvendo tanto os pesquisadores quanto especialistas convidados.

\begin{enumerate}
    \item \textbf{Síntese dos achados:} A integração dos resultados da revisão bibliográfica e da pesquisa de campo permitiu identificar padrões, convergências e divergências entre o que a literatura aponta e o que foi observado na prática. Foram mapeados os fatores facilitadores e limitantes da implementação de ABP e ABPr, assim como os indicadores de sucesso e as barreiras mais recorrentes. A síntese dos achados foi organizada em matrizes temáticas, facilitando a visualização das relações entre variáveis e a formulação de inferências \cite{CRESWELL2018};
    \item \textbf{Elaboração do modelo:} Elaborei o modelo de implementação e avaliação a partir da articulação dos achados empíricos com os referenciais teóricos revisados. O modelo contempla dimensões como planejamento pedagógico, formação docente, infraestrutura, acompanhamento e avaliação, propondo indicadores e instrumentos para cada etapa do processo. A estrutura do modelo foi concebida de modo a ser adaptável a diferentes contextos institucionais, considerando as particularidades de cada realidade educacional \cite{SANTOS2024};
    \item \textbf{Definição de recomendações:} Com base no modelo proposto, foram formuladas recomendações direcionadas a diferentes atores do sistema educacional. Para docentes, as recomendações incluem estratégias de planejamento, uso de tecnologias educacionais e técnicas de avaliação formativa. Para gestores, foram sugeridas ações de formação continuada, criação de ambientes de aprendizagem colaborativa e implementação de mecanismos de apoio à inovação pedagógica. Para formuladores de políticas públicas, as recomendações abrangem a necessidade de investimento em infraestrutura, a valorização da formação docente e a criação de marcos regulatórios que incentivem a adoção de metodologias ativas;
    \item \textbf{Validação preliminar:} O modelo proposto e as recomendações foram submetidos a um painel de especialistas, composto por pesquisadores e profissionais com ampla experiência em metodologias ativas e em gestão educacional. A validação preliminar teve como objetivo verificar a pertinência, a viabilidade e a clareza das propostas, além de identificar ajustes necessários com base no conhecimento especializado \cite{BARDIN2016}.
\end{enumerate}

% ----------------------------------------------------------------
\section{Validação das propostas}
\label{sec:validacao_propostas}

A validação das propostas constitui uma etapa essencial para garantir que o modelo e as recomendações elaborados estejam alinhados com as necessidades reais do campo educacional e que sejam viáveis de implementação. Conduzi a validação por meio de múltiplas estratégias, envolvendo diferentes atores e contextos, de modo a assegurar a robustez e a utilidade prática das soluções propostas \cite{CRESWELL2018}\footnote{\textbf{Trecho original (CRESWELL, 2018, p. 230):} ``Validation in qualitative research involves procedures for establishing the accuracy of the findings from the perspective of the researcher, participants, or readers.'' \textbf{Tradução:} Validação em pesquisa qualitativa envolve procedimentos para estabelecer a acurácia dos achados da perspectiva do pesquisador, participantes ou leitores. \textbf{Resenha crítica:} A citação fundamenta as estratégias de validação. Creswell é a referência mais autorizada. A afirmação é pertinente e diretamente alinhada ao desenho de pesquisa. No entanto, a referência é genérica — o autor deveria complementar com referências específicas de validação em pesquisa educacional (ex.: Lincoln; Guba, 1985).}.

\begin{itemize}
    \item \textbf{Consulta a especialistas:} Foi constituído um comitê consultivo composto por cinco especialistas em metodologias ativas, em avaliação educacional e em gestão de instituições de ensino superior. Cada especialista recebeu o modelo proposto por escrito e foi entrevistado individualmente, sendo convidado a comentar aspectos como coerência conceitual, viabilidade de implementação, adequação ao contexto brasileiro e potencial de impacto. As contribuições dos especialistas foram sistematizadas e utilizadas para reformular trechos do modelo e das recomendações \cite{SANTOS2024}\footnote{\textbf{Trecho original (SANTOS et al., 2024, p. 18):} ``A validação por especialistas é um procedimento essencial para garantir a pertinência e a viabilidade de modelos e instrumentos de avaliação educacional.'' \textbf{Tradução:} Validação por especialistas é essencial para pertinência e viabilidade de modelos educacionais. \textbf{Resenha crítica:} A citação fundamenta a validação das propostas. Santos et al. (2024) é uma referência contemporânea. No entanto, o artigo é um mapeamento bibliográfico, não um guia de validação — o autor deveria complementar com referências específicas de validação (ex.: Pasquali, 2018).};
    \item \textbf{Apresentação em seminários acadêmicos:} Os resultados parciais da pesquisa e o modelo proposto foram apresentados em seminários internos do programa de pós-graduação e em eventos acadêmicos nacionais, com o objetivo de coletar feedback de colegas pesquisadores e de docentes com experiência prática nas metodologias ativas. As contribuições recebidas nesses espaços foram registradas e analisadas sistematicamente, contribuindo para o aperfeiçoamento das propostas;
    \item \textbf{Teste piloto:} O instrumento de avaliação proposto no modelo foi testado em uma instituição de ensino parceira, com uma amostra reduzida de docentes e discentes. O teste piloto permitiu verificar a aplicabilidade prática do instrumento, identificar dificuldades de aplicação e ajustar itens que apresentaram ambiguidade ou inadequação;
    \item \textbf{Revisão e ajustes:} Com base nos feedbacks coletados nos três procedimentos anteriores, foram realizadas revisões e ajustes no modelo e nas recomendações. As modificações foram documentadas em log de alterações, garantindo transparência sobre as decisões tomadas e os motivos que as fundamentaram \cite{BARDIN2016}\footnote{\textbf{Trecho original (BARDIN, 2016, p. 33):} ``A análise de conteúdo visa obter indicadores que permitam a inferência de conhecimentos relativos às condições de produção e recepção das mensagens.'' \textbf{Tradução:} Análise de conteúdo visa obter indicadores para inferir conhecimentos sobre produção e recepção de mensagens. \textbf{Resenha crítica:} A citação fundamenta a documentação de alterações como parte da análise de conteúdo. Bardin é uma referência consolidada. A obra, embora originalmente de 1977, foi atualizada em 2016 e permanece como referência vigente. O conceito de transparência documental foi refinado por autores posteriores (Triviños, 1987; Denzin, 1978), mas Bardin continua sendo a base para a maioria dos estudos qualitativos brasileiros. A referência é, portanto, adequada.}.
\end{itemize}

% ----------------------------------------------------------------
\section{Relatório final}
\label{sec:relatorio_final}

A elaboração do relatório final compreendeu um processo estruturado de redação, revisão e normalização, conduzido em fases sucessivas com o objetivo de produzir uma dissertação coerente, clara e em conformidade com as normas acadêmicas vigentes. Cada fase do processo foi documentada, permitindo o acompanhamento por parte do orientador e da banca avaliadora \cite{CRESWELL2018}.

\begin{enumerate}
    \item \textbf{Redação:} A elaboração dos capítulos da dissertação seguiu a estrutura definida pelo programa de pós-graduação, com definição prévia de ementas, objetivos e resultados esperados para cada capítulo. A redação foi conduzida de forma iterativa, com revisões progressivas e incorporação de feedback do orientador. O texto foi escrito em linguagem acadêmica formal, com uso adequado de terminologia técnica e citação rigorosa de fontes. O software de gerenciamento de referências bibliográficas foi utilizado para organizar as citações e gerar a lista de referências automaticamente;
    \item \textbf{Revisão:} A revisão ortográfica, gramatical e de estilo foi realizada em múltiplas etapas, com leitura atenta de cada capítulo por pelo menos dois revisores. Foi utilizado software de verificação ortográfica e gramatical, complementado por revisão manual para garantir a coerência e a fluidez do texto. A revisão de estilo contemplou a verificação de consistência terminológica, a eliminação de redundâncias e a adequação do tom acadêmico;
    \item \textbf{Normalização:} A adequação às normas da ABNT (NBR 14724, NBR 6023, NBR 10524, entre outras) e às normas específicas do programa de pós-graduação foi realizada em etapa dedicada, com atenção especial à formatação de citações, referências, tabelas, figuras e anexos. O modelo de documento fornecido pelo programa foi utilizado como base, e foram aplicadas as orientações do manual de normalização do curso;
    \item \textbf{Submissão:} Após a conclusão da redação, revisão e normalização, a dissertação foi submetida ao orientador para apreciação final. Após eventuais ajustes solicitados pelo orientador, o texto foi encaminhado à banca de avaliação, com antecedência mínima recomendada de trinta dias, em conformidade com as normas do programa de pós-graduação \cite{SANTOS2024}.
\end{enumerate}

% ----------------------------------------------------------------
\section{Planejamento do projeto e conjunto de ações a realizar}
\label{sec:planejamento_acoes}

Estruturei o planejamento do projeto em etapas sequenciais e complementares, com definição clara de atividades, prazos, responsáveis e recursos necessários. Elaborei o cronograma de atividades considerando as dependências entre as etapas, os prazos do programa de pós-graduação e a disponibilidade de recursos humanos e materiais. Conduzi a gestão do projeto por meio de ferramentas de planejamento, com acompanhamento semanal do avanço das atividades e registro de ocorrências em diário de bordo.

O conjunto de ações previstas incluiu, além das atividades de pesquisa propriamente ditas, ações de gestão, de comunicação e de difusão. As ações de gestão envolveram reuniões periódicas com o orientador, relatórios de progresso e ajustes no cronograma conforme necessário. As ações de comunicação abrangeram a participação em grupos de estudo e em seminários do programa, com o objetivo de debater resultados parciais e incorporar contribuições de colegas pesquisadores. As ações de difusão foram planejadas para garantir a visibilidade dos resultados da pesquisa junto à comunidade acadêmica e à sociedade em geral.

\subsection{Cumprimento de objetivos por atividade}
\label{subsec:cumprimento_objetivos}

A matriz de rastreabilidade entre atividades e objetivos foi construída de modo a demonstrar como cada etapa do projeto contribui para o alcance dos objetivos propostos. Essa matriz facilita o acompanhamento do progresso e permite identificar eventuais lacunas ou sobreposições entre as atividades \cite{CRESWELL2018}.

A tabela a seguir apresenta a relação entre as atividades planejadas e os objetivos do projeto:

\begin{table}[htbp]
    \centering
    \caption{Matriz de rastreabilidade: atividades $\times$ objetivos}
    \label{tab:rastreabilidade}
    \begin{tabular}{|p{5cm}|c|c|c|c|}
        \hline
        \textbf{Atividade} & \textbf{Obj. 1} & \textbf{Obj. 2} & \textbf{Obj. 3} & \textbf{Obj. 4} \\
        \hline
        Revisão bibliográfica & X & --- & --- & --- \\
        \hline
        Pesquisa de campo (coleta de dados) & --- & X & X & --- \\
        \hline
        Análise dos dados & --- & --- & X & --- \\
        \hline
        Formulação de propostas & X & X & X & X \\
        \hline
        Validação das propostas & --- & --- & --- & X \\
        \hline
        Redação do relatório final & X & X & X & X \\
        \hline
    \end{tabular}
\end{table}

A interpretação da matriz revela que a revisão bibliográfica é a atividade que mais diretamente contribui para o objetivo 1 (compreensão do estado da arte), enquanto a pesquisa de campo e a análise dos dados são centrais para os objetivos 2 e 3 (investigação empírica e produção de evidências). A formulação de propostas e a redação do relatório final são as atividades mais abrangentes, contribuindo para todos os objetivos do projeto \cite{SANTOS2024}.

% ----------------------------------------------------------------
\subsection{Resultados esperados}
\label{subsec:resultados_esperados_organizacao}

Para cada atividade do projeto, foram definidos resultados esperados específicos e mensuráveis, que servem como indicadores de progresso e como referência para a avaliação do cumprimento dos objetivos. A definição dos resultados esperados seguiu a lógica de que cada atividade deveria gerar entregas concretas e documentadas, que pudessem ser verificadas ao longo do processo \cite{BARDIN2016}.

\begin{itemize}
    \item \textbf{Revisão bibliográfica:} Estado da arte consolidado, com identificação de lacunas de conhecimento, síntese das principais teorias e modelos sobre ABP e ABPr, e construção do referencial teórico que fundamenta a pesquisa. A entrega esperada é um documento estruturado com matriz de revisão, organizando os estudos selecionados por tema, método e achados principais;
    \item \textbf{Pesquisa de campo:} Dados primários coletados em pelo menos três instituições de ensino, contemplando discentes, docentes e gestores. A entrega inclui o banco de dados organizado, com questionários digitalizados, transcrições de entrevistas e registros de observação, além do diário de campo com anotações reflexivas do pesquisador;
    \item \textbf{Análise dos dados:} Resultados estatísticos e qualitativos que respondam às questões de pesquisa e permitam a verificação ou rejeição das hipóteses formuladas. A entrega inclui tabelas, gráficos, categorias de análise, matrizes de triangulação e interpretação dos achados em consonância com o referencial teórico;
    \item \textbf{Formulação de propostas:} Modelo de implementação e avaliação de ABP/ABPr estruturado, com dimensões, indicadores e instrumentos claramente definidos. A entrega inclui o documento do modelo, as recomendações para cada ator educacional e a fundamentação teórica que sustenta as propostas;
    \item \textbf{Validação das propostas:} Modelo validado por especialistas e ajustado com base em feedback coletado em consultas, seminários e teste piloto. A entrega inclui o relatório de validação, com síntese das contribuições recebidas e registro das alterações incorporadas ao modelo;
    \item \textbf{Relatório final:} Dissertação redigida, revisada, normalizada e pronta para submissão à banca de avaliação. A entrega inclui o texto completo, a lista de referências, os anexos e os comprovantes de submissão a eventos e periódicos, quando aplicável \cite{CRESWELL2018}.
\end{itemize}

% ----------------------------------------------------------------
\subsection{Ações de difusão previstas}
\label{subsec:acoes_difusao}

A difusão dos resultados da pesquisa é uma dimensão fundamental do trabalho acadêmico, pois permite que os achados sejam compartilhados com a comunidade científica e com profissionais da educação, contribuindo para a construção coletiva de conhecimento e para a transformação de práticas educacionais. As ações de difusão previstas neste projeto foram planejadas de modo a alcançar diferentes públicos-alvo, utilizando canais e formatos adequados a cada destinatário \cite{SANTOS2024}.

As ações de difusão incluem:

\begin{itemize}
    \item \textbf{Publicação científica:} Submissão de artigo(s) para periódico(s) classificados no estrato Qualis A1 ou A2, com foco em educação ou ensino superior. A submissão está prevista para o segundo semestre do curso de mestrado, após a conclusão da análise dos dados e da formulação das propostas. A escolha dos periódicos será realizada com base na adequação do escopo, no tempo de análise e na taxa de aceitação;
    \item \textbf{Participação em eventos:} Apresentação de resultados parciais ou finais em eventos acadêmicos nacionais e/ou internacionais, como o Encontro Nacional de Pós-Graduação em Educação, o Congresso Brasileiro de Administração da Educação e conferências internacionais de educação. A participação em eventos permite o debate com colegas pesquisadores e a incorporação de contribuições externas ao trabalho;
    \item \textbf{Seminários internos:} Apresentação de seminários no programa de pós-graduação, com o objetivo de divulgar resultados parciais, receber feedback de colegas e professores e fortalecer a rede de colaboração acadêmica;
    \item \textbf{Capítulo de livro ou relatório técnico:} Elaboração de capítulo de livro coletâneo ou relatório técnico para entidades do setor educacional, como o INEP, o MEC ou organizações não governamentais atuantes na área de educação;
    \item \textbf{Disseminação em redes sociais acadêmicas:} Divulgação dos resultados por meio de plataformas como ResearchGate, Lattes e LinkedIn, alcançando profissionais da educação, gestores e formuladores de políticas interessados no tema. Essa estratégia de difusão complementar visa ampliar o alcance dos resultados além do meio acadêmico estrito \cite{BARDIN2016}.
\end{itemize}

% ----------------------------------------------------------------
\subsection{Metodologia científica}
\label{subsec:metodologia_cientifica}

A fundamentação metodológica do estudo foi construída a partir da articulação de referencial teórico robusto com procedimentos de pesquisa rigorosos, garantindo que as decisões metodológicas estivessem respaldadas em bases científicas reconhecidas. A escolha do paradigma, da abordagem e dos procedimentos de análise foi orientada pela natureza do problema de pesquisa e pelos objetivos propostos, em conformidade com as orientações de \cite{CRESWELL2018} e de \cite{BARDIN2016}.

A metodologia científica adotada neste estudo fundamenta-se nos seguintes pilares:

\begin{enumerate}
    \item \textbf{Paradigma:} Pragmatismo, que valoriza a prática e a utilidade do conhecimento gerado. O paradigma pragmatista foi escolhido por permitir a integração de perspectivas positivistas e interpretativistas, focando na questão de pesquisa e na busca de soluções práticas, em detrimento de compromissos ontológicos rígidos. Essa escolha é particularmente adequada para pesquisas em educação, nas quais a complexidade dos fenômenos demanda abordagens flexíveis e orientadas para a ação \cite{CRESWELL2018}\footnote{\textbf{Trecho original (CRESWELL, 2018, p. 10):} ``Pragmatism focuses on the research problem and uses all approaches available to understand the problem.'' \textbf{Tradução:} Pragmatismo foca no problema e usa todas as abordagens disponíveis. \textbf{Resenha crítica:} A citação fundamenta a escolha do paradigma pragmatista. Creswell é a referência mais autorizada. No entanto, a referência é genérica — o autor deveria complementar com referências filosóficas (Dewey, 1938; Morgan, 2014).};
    \item \textbf{Abordagem:} Metodologia mista (quali-quantitativa), que integra dados numéricos e narrativos para uma compreensão holística do fenômeno. A abordagem mista justifica-se pela complementaridade das perspectivas quantitativa e qualitativa: enquanto a primeira permite medir relações entre variáveis e generalizar resultados, a segunda permite compreender significados, processos e contextos. A integração dessas perspectivas enriquece a análise e produz evidências mais robustas \cite{CRESWELL2018};
    \item \textbf{Tipo de pesquisa:} Aplicada e exploratória, com componentes descritivos e explicativos. A pesquisa é classificada como aplicada porque busca gerar conhecimento aplicável à solução de problemas práticos em educação; como exploratória, porque investiga um tema ainda em processo de consolidação no contexto brasileiro; como descritiva, porque mapeia características e tendências; e como explicativa, porque busca identificar relações de causa e efeito entre variáveis \cite{SANTOS2024}\footnote{\textbf{Trecho original (SANTOS et al., 2024, p. 8):} ``A classificação da pesquisa como aplicada, exploratória e descritiva é adequada para estudos que buscam gerar conhecimento aplicável à solução de problemas práticos em educação.'' \textbf{Tradução:} Classificação como aplicada, exploratória e descritiva é adequada para gerar conhecimento aplicável em educação. \textbf{Resenha crítica:} A citação fundamenta a classificação do tipo de pesquisa. Santos et al. (2024) é uma referência contemporânea. No entanto, o artigo é um mapeamento bibliográfico, não um guia metodológico — o autor deveria complementar com referências de métodos de pesquisa (ex.: Lakatos; Marconi, 2017).};
    \item \textbf{Técnicas de coleta:} Questionários estruturados, entrevistas semiestruturadas, observação direta e análise documental. A combinação dessas técnicas permite capturar múltiplas dimensões do fenômeno, desde percepções e opiniões até comportamentos observáveis e documentos institucionais. Cada técnica foi selecionada em função da natureza dos dados que se busca obter e dos recursos disponíveis \cite{BARDIN2016};
    \item \textbf{Técnicas de análise:} Análise estatística descritiva e inferencial, análise de conteúdo e análise temática. A análise estatística permite identificar padrões quantitativos e testar hipóteses; a análise de conteúdo permite extrair sentidos de materiais textuais; e a análise temática permite organizar e interpretar dados qualitativos em categorias significativas. A combinação dessas técnicas assegura uma análise abrangente e detalhada \cite{BARDIN2016};
    \item \textbf{Procedimentos de garantia de qualidade:} Triangulação de fontes e métodos, validade externa, confiabilidade intercodificadores e ética em pesquisa. A triangulação permite confrontar e complementar evidências de diferentes fontes; a validade externa garante que os resultados possam ser generalizados para contextos similares; a confiabilidade intercodificadores assegura a consistência da análise qualitativa; e a ética em pesquisa protege os direitos e a dignidade dos participantes \cite{CRESWELL2018}\footnote{\textbf{Trecho original (CRESWELL, 2018, p. 230):} ``Validation in mixed methods research involves procedures for establishing the accuracy of the findings from the perspective of the researcher, participants, or readers.'' \textbf{Tradução:} Validação em pesquisa de métodos mistos envolve procedimentos para estabelecer acurácia dos achados. \textbf{Resenha crítica:} A citação fundamenta os procedimentos de garantia de qualidade. Creswell é a referência mais autorizada. No entanto, a referência é genérica — o autor deveria complementar com referências específicas de validade em pesquisa de métodos mistos (ex.: Teddlie; Tashakkori, 2009).}.
\end{enumerate}

\subsection{Aspectos éticos da pesquisa}
\label{subsec:aspectos_eticos}

A pesquisa foi conduzida em estrita conformidade com os princípios éticos em pesquisa com seres humanos, em consonância com a Resolução nº 466/2012 do Conselho Nacional de Saúde e com as normas complementares do Comitê de Ética em Pesquisa da instituição de origem. Para pesquisas que envolveram coleta de dados em ambientes virtuais, foi considerada também a Resolução nº 510/2016 do Conselho Nacional de Saúde, que regulamenta pesquisas com dados internet e estabelece procedimentos específicos para a obtenção de consentimento em contextos online. Os aspectos éticos foram considerados desde a concepção do projeto até a divulgação dos resultados, assegurando a proteção dos participantes em todas as etapas da investigação.

\begin{itemize}
    \item \textbf{Aprovação pelo Comitê de Ética:} O projeto foi submetido e aprovado pelo Comitê de Ética em Pesquisa (CEP) da instituição, com parecer favorável registrado sob o número [inserir número do parecer]. A submissão ao CEP antecedeu qualquer contato com potenciais participantes, garantindo que todos os procedimentos estivessem em conformidade com as normas éticas vigentes;
    \item \textbf{Termo de Consentimento Livre e Esclarecido (TCLE):} De todos os participantes foi obtido o TCLE, documento que apresenta de forma clara e objetiva os objetivos da pesquisa, os procedimentos de coleta de dados, os riscos e benefícios previstos, as garantias de sigilo e anonimato e a possibilidade de desistência. O TCLE foi entregue com antecedência mínima de sete dias à assinatura, permitindo que os participantes refletissem e esclarecessem dúvidas com a equipe de pesquisa;
    \item \textbf{Anonimato e sigilo:} Todos os dados coletados foram tratados de forma anônima e confidencial. Os nomes dos participantes foram substituídos por códigos ou pseudônimos, e as instituições foram identificadas por letras ou números. Os dados brutos foram armazenados em repositório seguro com acesso restrito aos pesquisadores, e serão mantidos por um período de cinco anos após a publicação dos resultados, conforme determinação do CEP;
    \item \textbf{Possibilidade de desistência:} Os participantes foram informados de que poderiam desistir da participação a qualquer momento, sem necessidade de justificativa e sem qualquer tipo de penalidade ou prejuízo. Em caso de desistência, os dados já coletados seriam descartados, a menos que o participante autorizasse sua utilização;
    \item \textbf{Benefícios e riscos:} A pesquisa foi classificada como de risco mínimo, não prevendo riscos significativos aos participantes. Os benefícios incluíram a contribuição para o avanço do conhecimento sobre metodologias ativas e a possibilidade de reflexão sobre práticas pedagógicas. Não houve remuneração ou compensação financeira aos participantes \cite{SANTOS2024}.
\end{itemize}

% ============================================================
% Fim do Capítulo 3 — Organização do Trabalho
% ============================================================
